% appB-category.tex
\chapter{Category-Theoretic Language}\label{app:category}

The main text uses the vocabulary of categories sparingly, for three purposes: to say that a construction is universal, to describe quotients and products of representations, and to organize the order structure of representations. This appendix defines the terms and proves the statements used. No result here goes beyond finite sets and relations.

\section{Categories, functors, and the categories used here}

A \emph{category} has objects, morphisms between them, composition that is associative, and identities. A \emph{preorder} is a category with at most one morphism between any two objects, and a \emph{poset} is a preorder that is antisymmetric. The categories used in this volume are:
\begin{itemize}[nosep]
  \item $\mathbf{Set}$, with functions;
  \item the poset $\mathrm{Part}(X)$ of equivalence relations on a set $X$ ordered by inclusion (finer $\subseteq$ coarser);
  \item $\mathbf{Rep}(X)$, whose objects are surjections $r:X\to R$ and whose morphisms $r\to r'$ are functions $g:R\to R'$ with $r'=g\circ r$.
\end{itemize}

\begin{newresult}[Representations are partitions]\label{pb:rep}
There is at most one morphism between two objects of $\mathbf{Rep}(X)$, and a morphism $r\to r'$ exists iff $\ker r\subseteq\ker r'$. Two objects are isomorphic iff they have the same kernel. Hence $\mathbf{Rep}(X)$ is equivalent to the poset $\mathrm{Part}(X)$.
\end{newresult}
\assumes{None.}
\begin{proof}
If $g,g'$ are morphisms $r\to r'$, then $g(r(x))=r'(x)=g'(r(x))$ for all $x$, and $r$ is surjective, so $g=g'$. Existence is the factorization result (Result~\ref{thm:factorization}). Isomorphic objects have kernels contained in each other, hence equal; conversely equal kernels give mutually inverse factorizations. The assignment $r\mapsto\ker r$ is therefore an equivalence of categories onto $\mathrm{Part}(X)$.
\end{proof}

In this sense ``a representation'' and ``a partition of the cases'' are the same notion, which is the identification made in Chapter~\ref{ch:quotients}.

\section{Universal constructions}

\paragraph{Quotients as coequalizers.} Let $R\subseteq X\times X$ be a relation with projections $p_1,p_2:R\to X$. The \emph{coequalizer} of $p_1,p_2$ is a function $q:X\to Q$ with $q\circ p_1=q\circ p_2$ that is universal: every $f:X\to B$ with $f\circ p_1=f\circ p_2$ factors uniquely through $q$.

\begin{newresult}[Quotient by a relation]\label{pb:coeq}
In $\mathbf{Set}$ the coequalizer of $p_1,p_2:R\to X$ is the quotient map $q:X\to X/{\sim_R}$, where $\sim_R$ is the equivalence relation generated by $R$ (reflexive, symmetric, transitive closure).
\end{newresult}
\assumes{None.}
\begin{proof}
For $(x,y)\in R$ we have $q(x)=q(y)$, so $q\circ p_1=q\circ p_2$. If $f$ satisfies $f\circ p_1=f\circ p_2$ then $f(x)=f(y)$ for $(x,y)\in R$, and the set $\{(x,y):f(x)=f(y)\}$ is an equivalence relation containing $R$, hence containing $\sim_R$. So $f$ is constant on classes of $\sim_R$ and factors through $q$, uniquely since $q$ is surjective.
\end{proof}

\paragraph{Joint representations as products.} Given representations $r_1:X\to R_1$ and $r_2:X\to R_2$, the pairing $(r_1,r_2):X\to R_1\times R_2$ has kernel $\ker r_1\cap\ker r_2$.

\begin{newresult}[Meet and join of representations]\label{pb:meet-join}
In $\mathrm{Part}(X)$, ordered by inclusion of kernels (finer below coarser), the pair $\{\ker r_1,\ker r_2\}$ has a greatest lower bound $\ker r_1\cap\ker r_2$ (the joint representation) and a least upper bound $\ker r_1\vee\ker r_2$, the transitive closure of $\ker r_1\cup\ker r_2$ (the finest common coarsening). The joint representation is the categorical product in $\mathbf{Rep}(X)$ of $r_1$ and $r_2$: it maps to each of them, and any representation mapping to both factors through it. The finest common coarsening is the coproduct.
\end{newresult}
\assumes{None.}
\begin{proof}
The intersection of two equivalence relations is an equivalence relation, and it is the largest relation contained in both, hence the greatest lower bound. By Result~\ref{pb:rep} a representation $s$ maps to $r_i$ iff $\ker s\subseteq\ker r_i$, so $s$ maps to both iff $\ker s\subseteq\ker r_1\cap\ker r_2$, iff $s$ factors through the pairing. For the join, the transitive closure of the union is the smallest equivalence relation containing both, hence the least upper bound, and the universal property of the coproduct follows dually: $r_i$ maps to $s$ iff $\ker r_i\subseteq\ker s$, and both do iff $\ker r_1\cup\ker r_2\subseteq\ker s$ iff $\ker r_1\vee\ker r_2\subseteq\ker s$.
\end{proof}

The product is what Chapter~\ref{ch:adequacy} calls the joint representation for a family of questions; the coproduct is what the same chapter calls the common coarsening.

\section{Galois connections}

A \emph{Galois connection} between posets $P$ and $Q$ is a pair of order-reversing maps $F:P\to Q$, $G:Q\to P$ with $p\le G(F(p))$ and $q\le F(G(q))$. The composites $GF$ and $FG$ are closure operators.

\begin{newresult}[Questions and representations]\label{pb:galois}
Let $\mathcal F$ be a set of functions on $X$ (the questions), and let $\kappa(\Phi)=\bigcap_{f\in\Phi}\ker f$ for $\Phi\subseteq\mathcal F$. For an equivalence relation $\theta$ on $X$ let $\mathrm{Suff}(\theta)=\{f\in\mathcal F:\theta\subseteq\ker f\}$, the questions answerable from a representation with kernel $\theta$. Then $\Phi\mapsto\kappa(\Phi)$ and $\theta\mapsto\mathrm{Suff}(\theta)$ form a Galois connection between $(2^{\mathcal F},\subseteq)$ and $(\mathrm{Part}(X),\subseteq)$ (both order-reversing in the sense that larger sets of questions give finer kernels and finer kernels answer more questions), and $\Phi\subseteq\mathrm{Suff}(\kappa(\Phi))$. The closure $\mathrm{Suff}(\kappa(\Phi))$ is the set of all questions answerable by the least representation that answers $\Phi$.
\end{newresult}
\assumes{None.}
\begin{proof}
If $\Phi\subseteq\Psi$ then $\kappa(\Psi)\subseteq\kappa(\Phi)$, and if $\theta\subseteq\theta'$ then $\mathrm{Suff}(\theta')\subseteq\mathrm{Suff}(\theta)$; so both maps reverse inclusion. For $f\in\Phi$, $\kappa(\Phi)\subseteq\ker f$, hence $f\in\mathrm{Suff}(\kappa(\Phi))$. For an equivalence relation $\theta$, every $f\in\mathrm{Suff}(\theta)$ has $\theta\subseteq\ker f$, so $\theta\subseteq\kappa(\mathrm{Suff}(\theta))$. These are the two closure inequalities. By Result~\ref{pb:rep}, $f\in\mathrm{Suff}(\kappa(\Phi))$ iff $f$ factors through the representation with kernel $\kappa(\Phi)$.
\end{proof}

A consequence is that two families of questions are equivalent for the purposes of representation if they have the same closure, and that a representation designed for $\Phi$ automatically answers every other question in the closure at no extra cost.

\section{Functors and the history construction}

A \emph{functor} $F:\mathcal C\to\mathcal D$ assigns to objects and morphisms of $\mathcal C$ objects and morphisms of $\mathcal D$ preserving identities and composition. The assignment $\Hist\mapsto\{\Ext(h):h\in\Hist\}$ is not a functor on the nose; what is used in the text is the weaker statement that a simulation (a function between state sets commuting with transitions) induces a map of generated history sets. The reader who wants the full categorical story of minimal automata as terminal coalgebras will find it in Rutten, \emph{Universal coalgebra: a theory of systems}, and in Jacobs, \emph{Introduction to Coalgebra}; it is not needed here.

\section*{Exercises}
\addcontentsline{toc}{section}{Exercises}

\begin{exercise}
Describe $\mathbf{Rep}(X)$ for $|X|=3$ as a poset and draw its Hasse diagram. How many objects up to isomorphism?
\end{exercise}

\begin{exercise}
Compute the coequalizer in $\mathbf{Set}$ of the relation $R=\{(1,2),(3,4),(2,3)\}$ on $\{1,2,3,4,5\}$.
\end{exercise}

\begin{exercise}
For $X=\{0,\dots,11\}$, $r_1(x)=x\bmod3$, $r_2(x)=x\bmod4$, compute the kernel of the joint representation and of the finest common coarsening. How many classes does each have?
\end{exercise}

\begin{exercise}
Prove that in a Galois connection $F\circ G\circ F=F$.
\end{exercise}

\begin{exercise}
Show that for a family of questions $\Phi\subseteq\mathcal F$ the closure $\mathrm{Suff}(\kappa(\Phi))$ contains every question of $\mathcal F$ that is computable from the answers to the questions in $\Phi$, and give an example of a question in the closure that is not itself in $\Phi$.
\end{exercise}
